\documentclass[11pt,a4paper]{article}

\usepackage[utf8]{inputenc}
\usepackage[T1]{fontenc}
\usepackage{amsmath,amssymb,amsfonts,amsthm,bm}
\usepackage{graphicx}
\usepackage{booktabs}
\usepackage{hyperref}
\usepackage{url}
\usepackage[margin=1in]{geometry}
\usepackage{enumitem}

\newtheorem{theorem}{Theorem}
\newtheorem{corollary}[theorem]{Corollary}
\newtheorem{lemma}[theorem]{Lemma}

\title{Gauge Structure and Architectural Inductive Bias\\for Axial Parity Codes:\\A Theorem and a Negative Result}

\author{Brian Guarraci\\
  Galatheus\\
  \texttt{brian@galatheus.com}}

\date{May 2026}

\begin{document}
\maketitle

\begin{abstract}
We study the design of learned decoders for sparse-binary voxel
compression with an axial parity code. We prove that the codespace of
axial-line, plane, and global parity on the 3-torus is exactly the
image of the discrete third partial derivative $\partial_x \partial_y
\partial_z$ over $\mathbb{F}_2$, with closed-form rank $3n^2 - 3n + 1$,
identifying the parity codec as a discrete-$\mathbb{Z}_2$ instance of a
scalar-charge fracton gauge theory and placing it in the same algebraic
family as surface-code quantum error correction. The proof is elementary
once the lattice is read through its tensor-product structure, and the
theorem caught a real indexing bug in our previously deployed parity
matrix (reported rank $2016$ instead of the correct $2977$ at $n=32$).

The natural architectural reading of this theorem --- parameterise the
network's output as $x = \partial_x \partial_y \partial_z \phi$ and
predict the gauge potential $\phi$ --- fails decisively. Across $14$
configurations spanning capacity, training time, and several
optimisation tricks, this \emph{phi-prior} decoder plateaus at
validation IoU $\approx 0.29$ versus $\approx 0.48$ for an unconstrained
baseline at matched configuration; the gap is far larger than seed or
checkpoint-selection noise. A soft U($1$) latent-gauge variant does
better ($\approx 0.51$), but a controlled ablation shows its gauge
machinery is not load-bearing, and a generic multi-scale autoencoder
with no gauge structure of any kind wins outright ($0.572 \pm 0.004$
over three seeds). The principal lesson is methodological: a known
mathematical structure characterises \emph{what is reachable}, not the
right inductive bias for a learned decoder, and enforcing exact
algebraic structure architecturally can actively harm optimisation ---
a concrete coding-theoretic instance of the broader finding that hard
equivariance/constraint enforcement often underperforms a soft prior.
We report validation metrics on ShapeNet-$32^3$ and flag a held-out,
multi-seed cross-domain study (including surface-code QEC) as future
work.
\end{abstract}

\section{Introduction}
\label{sec:intro}

The parity codec for sparse-binary voxel compression encodes a $32^3$
binary occupancy field by transmitting (a) a syndrome consisting of
parity sums along axis-aligned lines, planes, and the volume total,
and (b) optional spectral side-information (low-frequency DCT
coefficients of the voxel field). At the receiver, a learned network
reconstructs the voxel field given the syndrome and side-information.
This is a clean instance of \emph{syndrome-conditioned reconstruction
on a regular lattice} and shares its abstract structure with
surface-code quantum error correction, where stabiliser measurements
play the role of parity checks.

Within the companion paper family, this note is the algebraic and
negative-control member. The compression paper asks which checks help a learned
proposal; the surface-code and X-cube papers test geometry-aware neural
decoders on quantum codes. Here we prove the algebraic structure behind the
axial-parity codec and then test the tempting architectural conclusion one
might draw from that proof. The main lesson is that a correct mathematical
description of the constraint set is not, by itself, a recipe for a good
learned decoder.

We pursue three connected goals:

\begin{enumerate}[nosep]
\item \textbf{Characterise the codespace algebraically.} What is the
structure of the set of voxel fields satisfying every parity check?
We give a theorem (\S\ref{sec:theorem}) stating that the codespace is
exactly the image of the discrete third partial derivative over
$\mathbb{F}_2$. This identifies the parity codec as a
discrete-$\mathbb{Z}_2$ instance of Pretko's scalar-charge fracton
gauge theory~\cite{pretko2017} and provides a constructive
parameterisation: every codeword is the discrete-XOR of eight corner
values of an underlying scalar potential $\phi$.

\item \textbf{Test whether the theorem prescribes a useful
architecture.} The most direct interpretation is to predict $\phi$
with a neural network and emit $x = \partial_x \partial_y \partial_z
\phi$ as a fixed output operator. We call this the \emph{phi-prior}
decoder. We find it fails decisively (\S\ref{sec:phi-prior}) ---
across capacity, training time, and several optimisation tricks, it
plateaus at $\approx 0.29$ IoU versus an unconstrained baseline's
$\approx 0.48$ at matched compute.

\item \textbf{Find what does work.} A multi-scale latent autoencoder
(\S\ref{sec:autoencoder}) --- encoder $\to 8^3$ latent $\to$ SharpNet
upsample, with no gauge structure of any kind --- beats both the
phi-prior and the unconstrained baseline. A controlled ablation
isolates the cause: the gauge machinery (both the hard $\mathbb{F}_2$
operator and the soft U($1$) regularisers) is not the active
ingredient; multi-scale processing is.
\end{enumerate}

The intended contribution is therefore a theorem, a negative result,
and a controlled ablation, organised around the question of when a
known mathematical structure should and should not be enforced
architecturally in a learned decoder. A natural follow-up --- whether
the same architectural lesson transfers across the algebraic family the
theorem unifies (in particular to surface-code QEC) --- is discussed as
a deferred, multi-seed cross-domain study in \S\ref{sec:qec}.

\section{Background}

\subsection{Axial parity codec}
Let $\Lambda = (\mathbb{Z}/n)^3$ be a periodic 3-cubic lattice with
voxel field $x : \Lambda \to \mathbb{F}_2$. The codec defines three
families of parity checks on $x$:
\begin{itemize}[nosep]
\item \textbf{Line.} $L^a_{\alpha\beta}(x) = \sum_t x_{t \hat{e}_a +
\alpha \hat{e}_b + \beta \hat{e}_c} \pmod 2$ for each axis $a \in
\{x,y,z\}$ and perpendicular coordinates $(\alpha, \beta)$.
\item \textbf{Plane.} $P^a_\alpha(x)$ summing over the plane perpendicular to $a$ at coord $\alpha$.
\item \textbf{Global.} $G(x) = \sum_v x_v$.
\end{itemize}
These checks are stacked into a parity matrix $H : \mathbb{F}_2^{n^3}
\to \mathbb{F}_2^M$. For $n=32$ the bug-fixed implementation has rank
$2977$. The codespace is $\mathcal{C}(n) = \ker H$.

The encoder transmits the syndrome $H x$ together with a low-frequency
DCT field of $x$ (variable bit budget; we use levels $k \in \{4,5,6,7\}$
giving $\{2744, 5832, 10648, 17576\}$ bits respectively). The decoder
constructs a canonical base solution $x_{\text{base}}$ from the
syndrome and a learned residual $r = x \oplus x_{\text{base}}$; the
residual lies in the codespace by construction since $x$ and
$x_{\text{base}}$ share the same syndrome. The learning task is to
predict $r$ from $x_{\text{base}}$, the syndrome, and (optionally) the
DCT side-info.

\subsection{Surface-code QEC}
The surface code is a $\mathbb{Z}_2$ stabiliser code on a $2$D toric
lattice~\cite{kitaev2003}; its decoder must infer the logical-error
class given the syndrome of stabiliser measurements (a $3$D detector
grid when stacked across rounds). It is a $\mathbb{Z}_2$ $1$-form
gauge theory in the Wegner--Kitaev sense~\cite{wegner1971,kitaev2003},
one rung below the fracton gauge theory the parity codec instantiates.
Our deployed surface-code decoder is a flat-resolution $3$D
ResNet~\cite{he2016}.

\subsection{Notation}
$(\partial_a \phi)_v := \phi_v + \phi_{v + \hat e_a} \pmod 2$ is the
discrete forward difference over $\mathbb{F}_2$ in direction $a$.

\section{Theorem: codespace structure of the axial parity code}
\label{sec:theorem}

\begin{theorem}\label{thm:codespace}
For all $n \geq 1$ on the $3$-torus $\Lambda = (\mathbb{Z}/n)^3$,
\[
\mathcal{C}(n) = \mathrm{image}\bigl(\partial_x \partial_y \partial_z\bigr)
\subseteq \mathbb{F}_2^{n^3}, \qquad
\mathrm{rank}_{\mathbb{F}_2}(H) = 3n^2 - 3n + 1.
\]
\end{theorem}

\begin{proof}
Write $V = \mathbb{F}_2^n$ and identify
$\mathbb{F}_2^{n^3} = V \otimes V \otimes V$, the three factors indexed
by the $x,y,z$ coordinates. The one-direction forward difference acts on
a single tensor factor: $\partial_x = \partial \otimes I \otimes I$,
$\partial_y = I \otimes \partial \otimes I$, $\partial_z = I \otimes I
\otimes \partial$, where $\partial = I + S$ and $S$ is the cyclic shift
on $V$. Hence $T := \partial_x \partial_y \partial_z = \partial \otimes
\partial \otimes \partial$.

\smallskip
\noindent\emph{Step 1 (codespace as an orthogonal complement).} Each
line, plane, and global check equals the $\mathbb{F}_2$ inner product of
$x$ with the indicator vector of an axis-parallel line, an axis-normal
plane, or the whole lattice, respectively. Thus
$\mathcal{C}(n) = \ker H = R^\perp$, where $R \subseteq
\mathbb{F}_2^{n^3}$ is the span of all check indicators. A plane is the
$\mathbb{F}_2$-sum of the parallel lines it contains, and the global
vector is the sum of parallel planes, so $R$ is already spanned by the
axis-parallel line indicators alone.

\smallskip
\noindent\emph{Step 2 (the span of lines).} The indicator of the
$x$-line at transverse coordinates $(\beta,\gamma)$ is $\mathbf{1}
\otimes \delta_\beta \otimes \delta_\gamma$, where $\mathbf{1} \in V$ is
the all-ones vector and $\delta_j$ are standard basis vectors. As
$(\beta,\gamma)$ range over $[n]^2$ the vectors $\delta_\beta \otimes
\delta_\gamma$ span $V \otimes V$, so the $x$-lines span $\langle
\mathbf{1} \rangle \otimes V \otimes V$. This is exactly $K_x := \{f :
\partial_x f = 0\}$, the fields constant along $x$. Likewise the $y$-
and $z$-lines span $K_y$ and $K_z$, so $R = K_x + K_y + K_z$.

\smallskip
\noindent\emph{Step 3 (dimension).} Choose any basis $B$ of $V$ with
$\mathbf{1} \in B$. Then $K_x, K_y, K_z$ are each spanned by a subset of
the common tensor basis $B \otimes B \otimes B$, hence are coordinate
subspaces, so inclusion--exclusion for dimensions is exact. With
$\dim K_a = n^2$, $\dim(K_a \cap K_b) = n$ (fields constant in two axes),
and $\dim(K_x \cap K_y \cap K_z) = 1$ (global constants),
\[
\dim R = 3n^2 - 3n + 1.
\]
Therefore $\mathrm{rank}_{\mathbb{F}_2}(H) = \dim R = 3n^2 - 3n + 1$ and
$\dim \mathcal{C}(n) = n^3 - (3n^2 - 3n + 1)$.

\smallskip
\noindent\emph{Step 4 (kernel of $T$).} For a tensor product of
operators acting on separate factors, $\ker(\partial \otimes \partial
\otimes \partial) = \ker\partial \otimes V \otimes V + V \otimes
\ker\partial \otimes V + V \otimes V \otimes \ker\partial$. On the
$n$-cycle, $\ker\partial = \langle \mathbf{1} \rangle$ (a vector killed
by $I+S$ is shift-invariant, hence constant). Thus $\ker T = K_x + K_y +
K_z = R$, and $\dim \mathrm{image}(T) = n^3 - \dim\ker T = n^3 - (3n^2 -
3n + 1) = \dim \mathcal{C}(n)$.

\smallskip
\noindent\emph{Step 5 (equality).} For any $\phi$, every axis-aligned
line/plane/global sum of $T\phi = \partial_x \partial_y \partial_z \phi$
telescopes to $0 \pmod 2$ under the periodic boundary, since each
$\phi_v$ enters such a sum an even number of times; hence
$\mathrm{image}(T) \subseteq \mathcal{C}(n)$. Two subspaces, one
contained in the other and of equal finite dimension, coincide.
Therefore $\mathcal{C}(n) = \mathrm{image}(T)$.
\end{proof}

\paragraph{Empirical verification.}
We verified the theorem by brute-force $\mathbb{F}_2$ Gaussian
elimination on the full check matrix at $n \in \{4, 6, 8\}$
(\texttt{verify\_fracton\_theorem.py}), obtaining
$\mathrm{rank}(H) \in \{37, 91, 169\}$ in exact agreement with
$3n^2 - 3n + 1$, $\mathrm{rank}(T) \in \{27, 125, 343\}$ in agreement
with $n^3 - (3n^2 - 3n + 1)$, and $\dim \mathcal{C}(n) = \mathrm{rank}(T)$
with random $\phi$ producing codewords that satisfy every check
($50/50$ per $n$), confirming $\mathcal{C}(n) = \mathrm{image}(T)$ as a
subspace. The $n=32$ arithmetic gives $3 \cdot 1024 - 96 + 1 = 2977$.

\paragraph{The bug the theorem caught.}
The previously deployed parity matrix reported rank $2016$ at $n=32$
rather than $2977$. Running its construction through the same verifier
yields $\mathrm{rank} \in \{28, 66, 120\}$ at $n \in \{4,6,8\}$ ---
strictly smaller than $\{37,91,169\}$, hence a \emph{different,
lower-rank} code (a duplicate-axis indexing error, not a
re-presentation of the same checks). The closed-form rank is thus also
a practical correctness check for the codec's bit-budget accounting.

\paragraph{Interpretation.}
The codespace is the discrete-$\mathbb{F}_2$ analogue of Pretko's
continuum scalar-charge fracton gauge theory~\cite{pretko2017}: the
codeword $x$ is the third-derivative ``field strength'' of an underlying
potential $\phi$. The conserved charges of the fracton matter sector ---
total charge ($G$), planar charges ($P^a_\alpha$), and line charges
($L^a_{\alpha\beta}$) --- are exactly the parity checks of the codec.
Surface-code QEC corresponds to one rung up the fracton hierarchy: a
$\mathbb{Z}_2$ $1$-form gauge theory whose Gauss law is the standard
four-weight stabiliser. The fracton$\leftrightarrow$code dictionary is
itself established~\cite{pretko2017,vijay2016,haah2011}; our contribution
is the specific identification of the axial-parity codec as the
$\partial_x\partial_y\partial_z$ field-strength theory and the
closed-form rank that follows.

\section{Experimental setup}
\label{sec:setup}

All experiments use ShapeNet-$32^3$~\cite{shapenet} binary voxels with
the bug-fixed parity check matrix giving syndrome of dimension $M =
2977$. Training target is the residual $r = x \oplus x_{\text{base}}$,
which lies in the codespace. Loss is binary cross-entropy on $r$,
optionally with positive-class weighting. The reported metric is
intersection-over-union (IoU) on the reconstructed voxel field $\hat x =
x_{\text{base}} \oplus \hat r$, evaluated on a held-back validation set
(not a disjoint test split; see \S\ref{sec:limitations}).

Unless stated otherwise: train subset $2{,}048$ samples, validation
$512$ samples, batch size $8$, AdamW with lr $5 \times 10^{-4}$, cosine
LR schedule. We report best validation IoU achieved during training,
and multi-seed mean $\pm$ stdev where indicated.

Architectures compared:
\begin{itemize}[nosep]
\item \textbf{x-prior} (baseline): ResNet3D at varying width/depth, predicts $r$ directly via flat-resolution $3$D convolution at constant spatial resolution, syndrome conditioning via MLP.
\item \textbf{phi-prior:} same backbone, output reinterpreted as $\phi$; final operator is the bilinear soft-XOR of $8$ corner $\phi$'s with output bias init at $-5$ to break XOR symmetry.
\item \textbf{gauge transporter:} U($1$) latent-lattice gauge with parallel-transport reconstruction from a seed voxel, plus auxiliary curvature/transport regularisers.
\item \textbf{latent autoencoder:} encoder $\to 8^3$ latent $\to$ SharpNet decoder, no gauge structure.
\end{itemize}

\section{Hard architectural enforcement plateaus far below baseline}
\label{sec:phi-prior}

Across $14$ experiments varying capacity, training time, and
optimisation tricks (output bias init, positive-class weighting,
residual escape hatch, straight-through estimator, multi-scale
capacity), the phi-prior decoder plateaus at val\_iou $\approx 0.29$
vs the unconstrained x-prior at $0.48$ at matched parameter count and
training budget.

\begin{table}[h]
\centering
\small
\begin{tabular}{lrrrr}
\toprule
variant & params & epochs & pos\_weight & val\_iou \\
\midrule
x-prior                                  & 5.0M  & 20 & 19 & \textbf{0.481} \\
x-prior                                  & 12.0M & 20 & 19 & 0.489 \\
x-prior                                  & 17.9M & 20 & 19 & 0.474 \\
\midrule
phi-prior, tanh-product XOR              & 1.7M  &  3 &  1 & 0.002 \\
phi-prior, bilinear + bias init $-5$     & 1.7M  &  5 &  1 & 0.021 \\
phi-prior, $+$pos\_weight$=19$           & 5.0M  & 20 & 19 & 0.283 \\
phi-prior, 12M                           & 12.0M & 20 & 19 & 0.276 \\
phi-prior, 30 epochs                     & 5.0M  & 30 & 19 & \textbf{0.291} \\
phi-prior, 8K train data                 & 5.0M  & 25 & 19 & 0.292 \\
phi-prior, residual escape hatch         & 5.0M  &  5 &  1 & 0.001 \\
phi-prior, STE hard-XOR                  & 1.7M  & 10 &  1 & 0.001 \\
\bottomrule
\end{tabular}
\caption{Phi-prior decoder vs x-prior baseline (best validation IoU,
single-seed per cell). The phi-prior plateau at $\approx 0.29$ IoU is
robust across capacity, training duration, and optimisation; the
unconstrained x-prior at matched compute reaches $\approx 0.48$. The
robustness of the gap across $14$ cells, not any single comparison, is
the evidence.}
\label{tab:phi-prior}
\end{table}

The plateau is robust to capacity ($5$M $=$ $12$M), training duration
($20=30$ epochs), positive weighting (sweet spot at pos\_weight$=19$),
and data scale ($2$K $=$ $8$K). The phi-prior occasionally exhibits a
phase-transition jump from $\sim 0.13$ to $\sim 0.21$ around epoch $12$
under the right hyperparameters but always saturates at $\sim 0.29$
soon after. \textbf{Architectural enforcement of the codespace is
therefore not a viable inductive bias on this task.}

The reason is mechanical: the bilinear soft-XOR over $8$ corners has
gradient bounded by $1$ but multiplicatively attenuated through the
chain, and the network must simultaneously modulate $8$ voxel-position
$\phi$ values to flip a single residual bit --- a non-local
credit-assignment problem that the BCE loss cannot drive efficiently
when the target distribution is sparse (most residuals are $0$).

\section{What works instead: a generic multi-scale prior}
\label{sec:autoencoder}

If hard gauge enforcement is the wrong inductive bias, what is the
right one? A generic multi-scale latent autoencoder --- encoder
downsampling $32^3 \to 8^3$ with $C_{\mathrm{latent}}$ channels, then a
SharpNet decoder upsampling $8^3 \to 32^3$ with residual blocks at each
scale, and \emph{no gauge structure of any kind} --- beats both the
phi-prior and the unconstrained x-prior at matched compute
(Table~\ref{tab:compare}).

\begin{table}[h]
\centering
\small
\begin{tabular}{llrr}
\toprule
architecture & gauge structure & params & val\_iou \\
\midrule
phi-prior (hard $\mathbb{F}_2$, every voxel) & hard                  & 5.0M  & 0.283 \\
x-prior (unconstrained ResNet3D)             & none                  & 5.0M  & 0.481 \\
soft U($1$) transporter                      & soft, coarse latent   & 4.4M  & 0.509 \\
\quad --- auxiliary gauge losses removed     & soft, coarse latent   & 4.4M  & 0.518 \\
\textbf{multi-scale autoencoder}             & \textbf{none}         & 16M   & \textbf{$0.572 \pm 0.004$} \\
\bottomrule
\end{tabular}
\caption{Architecture comparison (best validation IoU on $512$
held-back ShapeNet samples). The autoencoder row is the mean $\pm$
stdev over seeds $\{42,43,44\}$; the other rows are single-seed (the
phi-prior plateau is corroborated by the $14$-run sweep in
Table~\ref{tab:phi-prior}). Hard gauge enforcement is worst; soft gauge
structure helps modestly but its machinery is not load-bearing; the
gauge-free multi-scale autoencoder wins.}
\label{tab:compare}
\end{table}

Two controlled observations isolate the cause. First, the soft U($1$)
transporter with all auxiliary curvature/transport/latent-reconstruction
losses \emph{removed} (val\_iou $0.518$) slightly outperforms the full
version ($0.509$): the gauge regularisers are not load-bearing. Second,
dropping the parallel-transport operator entirely --- leaving a plain
multi-scale autoencoder --- improves things further ($4.1$M: $0.554$;
$16$M: $0.572$). The win is multi-scale processing, not gauge structure.

A single-seed bottleneck-size sweep indicates an $8^3$ latent is
near-optimal: a tighter $4^3$ bottleneck loses information ($0.562$ at
$32$M) and a looser $16^3$ underconstrains ($0.542$ at $3$M). This is
consistent with the interpretation that the bottleneck implements an
implicit low-frequency/smoothness prior, matched to the sparse,
surface-concentrated structure of the occupancy field --- precisely the
property the gauge constraint is orthogonal to. We treat the
bottleneck-sweep numbers as single-seed exploratory evidence, not
multi-seed claims.

\section{A deferred cross-domain question}
\label{sec:qec}

The gauge theorem places axial-parity compression, surface-code QEC, and
X-cube fracton QEC in a common constraint-system family: the surface code is a
$\mathbb{Z}_2$ $1$-form gauge theory, while the parity codec is a
discrete-$\mathbb{Z}_2$ scalar-charge fracton gauge theory. It is natural to ask
whether the multi-scale prior that wins here also helps for QEC syndrome
decoding, where the syndrome is a dense $3$D detector volume rather than a
sparse occupancy field.

Preliminary single-seed experiments suggest the opposite: aggressive
spatial bottlenecks appear to destroy the per-detector-cell information
that dense surface-code syndromes carry, where a deployed
flat-resolution convolution preserves it. We deliberately do \emph{not}
report these as a result. A defensible cross-domain claim requires a
disjoint test split, multiple seeds, and matched baselines decoded on
the same shots --- the standard we apply to the compression numbers
above is not yet met on the QEC side. We state only the hypothesis the
preliminary runs suggest: that the multi-scale prior's advantage is
specific to sparse signals with a smoothness structure and does not
extend to dense per-cell syndromes. Testing it rigorously is future
work.

\section{Discussion}

\paragraph{A theorem characterises feasibility, not inductive bias.}
The gauge theorem tells us \emph{what is reachable} (the codespace),
not \emph{how to reach it} (the architecture). Architectural
enforcement of the theorem's constructive parameterisation hurt; the
right inductive bias for this task is unrelated to the gauge structure
and is well-served by a generic multi-scale architecture. The
compression codec's existing DCT side-information implements a related
smoothness prior explicitly, and the multi-scale autoencoder appears to
learn one implicitly via its $8^3$ bottleneck.

\paragraph{Relation to soft-constraint priors.}
The phi-prior failure is a concrete coding-theoretic instance of a more
general phenomenon: hard architectural enforcement of an exact
algebraic or equivariance structure can underperform a soft prior or an
unconstrained pathway. Finzi et al.'s residual pathway
priors~\cite{finzi2021rpp} make this point for equivariance constraints
and propose converting hard constraints into soft penalties; our soft
U($1$) transporter is the analogous move in this setting, and the
ablation shows even its soft gauge losses are not what carries the win.
The lesson is not that the gauge structure is wrong --- it is exactly
correct as a characterisation of the codespace --- but that exactness
of the constraint and usefulness as an inductive bias are different
properties.

\paragraph{Open questions.}
(i)~Can the gauge structure be exploited as a \emph{post-hoc}
projection or a \emph{soft penalty} (rather than as architectural
enforcement) without the optimisation cost we observed?
(ii)~The theorem generalises to higher dimensions (the codespace of
axial parity on the $d$-torus is the image of $\partial_1 \cdots
\partial_d$, with rank obtainable by the same tensor argument); does the
architectural lesson generalise too? (iii)~Does the
sparse-with-smoothness vs dense-per-cell distinction sketched in
\S\ref{sec:qec} hold up under a rigorous multi-seed cross-domain test,
and where do quantum LDPC codes with sparser stabilisers fall?

\section{Limitations}
\label{sec:limitations}

We are explicit about scope. (1)~All reported IoU is \emph{validation}
IoU on a held-back ShapeNet-$32^3$ split, not a disjoint held-out test
split; the multi-seed autoencoder result uses seeds $\{42,43,44\}$,
while the x-prior, phi-prior, and transporter rows and the
bottleneck-size sweep are single-seed exploratory measurements.
(2)~The matched-compute comparison in Table~\ref{tab:compare} is
matched in training budget but not exactly in parameter count across all
rows; the headline conclusion (hard enforcement worst, gauge machinery
not load-bearing, multi-scale wins) rests on gaps far larger than seed
noise and on the directionality of the two ablations rather than on a
single point estimate. (3)~Results are at $32^3$ resolution on a single
object dataset. (4)~The cross-domain QEC discussion
(\S\ref{sec:qec}) is a hypothesis, not a measured result.

\section{Conclusion}

We proved a clean algebraic characterisation of the parity codec's
codespace --- it is exactly the image of the discrete third derivative
$\partial_x \partial_y \partial_z$ over $\mathbb{F}_2$, with closed-form
rank $3n^2 - 3n + 1$ --- connecting it to fracton gauge theory and to
surface-code QEC, and we used it to catch a real indexing bug in a
deployed parity matrix. We then tested the natural architectural
translation of the theorem (predicting the gauge potential and applying
the third derivative as a fixed output operator) and found it
catastrophically worse than an unconstrained baseline. The actual win on
the compression task comes from a generic multi-scale autoencoder with
no gauge structure, and a controlled ablation shows the gauge machinery
is not the active ingredient.

The principal lesson is methodological: a known mathematical structure
characterises feasibility but does not prescribe the right inductive
bias; architectural choice should be governed by the optimisation
landscape and signal structure of the specific task. We provide a
verified theorem, a clean negative example of constraint enforcement,
and the ablation that localises why it fails.

\paragraph{Reproducibility.}
All experiments are reproducible from the released code. In the family
repository, the theorem verifier is grouped with the X-cube fracton code, and
the gauge negative-result experiments are preserved as their own reproducible
submodule. Run-by-run
configuration and results, including discarded approaches, are
catalogued in \texttt{experiments.tsv}. The theorem verifier
(\texttt{verify\_fracton\_theorem.py}) reproduces the rank claims at
$n \in \{4,6,8\}$ in $\sim 20$~s on CPU. See the companion
\texttt{REPRODUCIBILITY.md} for exact commands.

% ====================================================================
\begin{thebibliography}{99}

\bibitem{pretko2017} M. Pretko.
\newblock Subdimensional particle structure of higher rank U(1) spin liquids.
\newblock {\em Phys.~Rev.~B}, 95:115139, 2017.

\bibitem{vijay2016} S. Vijay, J. Haah, and L. Fu.
\newblock Fracton topological order, generalized lattice gauge theory, and duality.
\newblock {\em Phys.~Rev.~B}, 94:235157, 2016.

\bibitem{haah2011} J. Haah.
\newblock Local stabilizer codes in three dimensions without string logical operators.
\newblock {\em Phys.~Rev.~A}, 83:042330, 2011.

\bibitem{wegner1971} F. Wegner.
\newblock Duality in generalized {I}sing models and phase transitions without local order parameters.
\newblock {\em J.~Math.~Phys.}, 12:2259, 1971.

\bibitem{kitaev2003} A.~Y. Kitaev.
\newblock Fault-tolerant quantum computation by anyons.
\newblock {\em Annals of Physics}, 303(1):2--30, 2003.
  arXiv:quant-ph/9707021 (1997).

\bibitem{shapenet} A.~X. Chang et al.
\newblock {ShapeNet: An Information-Rich 3D Model Repository}.
\newblock arXiv:1512.03012, 2015.

\bibitem{he2016} K. He, X. Zhang, S. Ren, and J. Sun.
\newblock Deep residual learning for image recognition.
\newblock In {\em CVPR}, 2016.

\bibitem{finzi2021rpp} M. Finzi, G. Benton, and A.~G. Wilson.
\newblock Residual pathway priors for soft equivariance constraints.
\newblock In {\em NeurIPS}, 2021. arXiv:2112.01388.

\end{thebibliography}

\end{document}
